%% ma: L12-B14-0027
%% lop: 12
%% chuong: 5
%% bai: 14
%% dang: 12.14.16
%% loai: TN4
%% muc: TH
%% dap_an: "C"
%% nguon: "CĐ2.5.pdf giáo viên gửi 27/10/2026 (M07, Lê Minh Khôi – Nguyễn Anh Khôi, Chương II Chủ đề 5 Tích có hướng của hai vectơ), Câu 11"
%% kien_thuc: "Bài 14 (tích có hướng, phương trình mặt phẳng qua ba điểm)"
%% trang_thai: chinh-thuc
%% ghi_chu: "Giáo viên đã duyệt 27/10/2026"
\begin{ex}
Trong không gian $Oxyz$, cho bốn điểm $A(1;1;4)$, $B(5;-1;3)$, $C(2;2;m)$ và $D(3;1;5)$ với $m\in\mathbb R$. Tìm giá trị của $m$ để bốn điểm trên tạo thành tứ diện.
\choice{$m\in\mathbb R$}{$m\ne4$}{\True $m\ne6$}{$m<0$}
\loigiai{$\left[\overrightarrow{AB},\overrightarrow{AD}\right]=(-2;-6;4)$ nên $(ABD)$: $x+3y-2z+4=0$.

$C\notin(ABD)\Leftrightarrow2+6-2m+4\ne0\Leftrightarrow m\ne6$.}
\end{ex}
